% State of the Art section. Paste into your own thesis, or \input it.
%
% Built from docs_thesis/literature.md Axis 2.
% Verification status: all entries below are [BIBLIOGRAPHY VERIFIED], abstract-level.

\section{Describing the tree}
\label{sec:describing-tree}

Once a tree can be built reliably (Section~\ref{sec:getting-tree}), the next question is what a
population of them looks like.
Two studies have described one, at a scale in the hundreds.

Medrano-Gracia et al.\ built a computational atlas of coronary bifurcation geometry over 300 CCTAs
selected specifically for a zero calcium score and no stenosis \cite{medranogracia2016atlas}.
Automatic quantification of 3D angles, diameters and lengths across the tree gave a left main
diameter of $3.5 \pm 0.8$~mm and length $10.5 \pm 5.3$~mm, and a left main distal bifurcation angle of
$89 \pm 21^\circ$ in patients with a ramus intermedius against $75 \pm 23^\circ$ without ($p<0.001$),
with equivalent tabulations for the LAD/diagonal, LCX/OM and right crux junctions.
Reporting a population as a spread of values rather than a single mean is the reporting shape
objective 8 needs, and the split by the presence of a ramus intermedius is an external result to
check the ImageCAS-X cohort's own IM-versus-no-IM angle contrast against
(Section~\ref{sec:dataset-dominance}).
Their cohort's selection matters wherever the numbers are compared to this thesis's: it is
healthy-selected, zero calcium and no stenosis, where ImageCAS-X is 388 of 800 cases with CAD by
\texttt{Descriptors.xlsx}, so the two populations are not directly comparable without saying so.
A companion paper carries the same cohort into more depth on the shape modeling itself
\cite{medranogracia2017bifurcation}, and both are distributed through the public Coronary Atlas
resource.

Nannini et al.\ automated the whole pipeline from image to geometry in 281 CCTA patients from a
single center: a cascaded 2.5D-then-3D U-Net segmentation (Dice 0.896), VMTK centerline extraction,
geometric characterization, and stratification by cardiovascular risk factor
\cite{nannini2024characterization}.
Three things are worth adopting directly.
Their tortuosity is defined three ways --- a local arc-length-to-chord ratio over 1~cm reference arcs,
a direction-change tortuosity angle, and the clinical count of bends past $45^\circ$ that Zebi\'c
Mihi\'c et al.\ also use \cite{zebicmihic2023tortuosity} (Section~\ref{sec:geometry}) --- which gives
objective 7 a published definition to match rather than an invented one.
Their zoning of a vessel into proximal, medial and distal thirds by normalized geodesic distance
from the ostium is the along-vessel parameterization objective 7 needs and does not yet have.
And their headline finding, a negative correlation between tortuosity and calcific plaque volume that
strengthens distally as calcium falls, is itself a population-scale geometry-disease association at
the vessel level, the kind Section~\ref{sec:geometry} otherwise had to draw from cohorts in the low
hundreds.
What their pipeline does not do is what Section~\ref{sec:the-gap} turns on: no bifurcation angles, no
dominance, no graph or tree model of the whole tree, and no attempt to flag which of the 281 patients
is unusual. Single-scanner training data and a 43-case test set are the authors' own stated limits.

Both studies stop at description.
Neither uses the distribution it built to ask which individual falls outside it.
